\section{XJTU-SY 结果}

\subsection{主划分设置与证据来源}

XJTU-SY 的主结果来自 Step F，即 \config{xjtu_bearing_index_split} + \config{manual_basic} + \config{degradation_three_tasks}。

该 run 的 train/val/test 都覆盖三种工况，但按轴承后缀隔离。

训练集只用于 ranking。

验证集和测试集用于分布检查、曲线检查和图像审查。

Step F 的输出包括：

\begin{itemize}[leftmargin=2em]
  \item 结果目录：\path{xjtu_sy/all_conditions_bearing_index_manual_basic/}
  \item 排名与建议：\path{feature_ranking.csv}、\path{feature_recommendations.md}
  \item RUL 与退化证据：\path{rul_correlation.csv}、\path{degradation_scores.csv}
  \item 分类证据：\path{health_state_separability.csv}、\path{early_fault_scores.csv}
\end{itemize}

Step H 和 Step I 不是新结论的替代品，而是稳定性检查。

Step H 回答同一工况内是否稳定。

Step I 回答跨工况是否仍可解释。

\begin{table}[htbp]
\centering
\caption{XJTU-SY main split 的任务级 top evidence}
\label{tab:xjtu_feature_evidence}
\scriptsize
\begin{tabularx}{\textwidth}{llrY}
\toprule
任务 & 特征 & score/rank & 解释 \\
\midrule
RUL & \feature{mag__time__mean} & rank 2, score 0.625 & Spearman=-0.768，趋势接近 label-source RMS，但不是实际 HI source。\\
RUL & \feature{mag__time__mean_abs} & rank 2, score 0.625 & 与 magnitude mean 同分，适合 compact baseline 二选一或一起做消融。\\
RUL & \feature{h__time__mean_abs} & rank 4, score 0.619 & 水平通道幅值候选；condition-wise top10 count=3。\\
Health State & \feature{h__time__mean_abs} & rank 1, score 0.996 & MI=0.623, Fisher=1.864, ANOVA F=4366.3；主 split boxplot 分离强。\\
Health State & \feature{mag__time__mean} & rank 4, score 0.875 & 跨工况后更稳，是最终 A 级推荐。\\
Early Fault & \feature{h__time__mean_abs} & rank 1, score 0.996 & 主 split 很强，但 Step H/I 后降级为 condition-sensitive。\\
Early Fault & \feature{mag__time__mean} & rank 4, score 0.961 & Step H top10 count=3，Step I 保持 stable cross-condition。\\
Reference & \feature{mag__time__rms} & RUL rank 1 & 实际 HI/FPT source，仅作 sanity reference。\\
\bottomrule
\end{tabularx}
\end{table}


\subsection{RUL 特征分析}

XJTU-SY 的 RUL top features 主要来自 magnitude 和 horizontal amplitude family。

主划分中 \feature{mag__time__rms} 的 RUL rank 为 1，score 为 0.625。

但它是实际 HI/FPT label-source，因此不能作为独立候选。

独立候选应优先看 \feature{mag__time__mean}、\feature{mag__time__mean_abs} 和 \feature{mag__time__std}。

\feature{mag__time__mean} 在 Step F 中 rank 2，RUL score 为 0.625。

它的 Pearson 为 -0.792，Spearman 为 -0.768，Kendall 为 -0.628。

负相关不是问题，因为 RUL 越小表示越接近失效，而振动幅值通常越大。

因此绝对 Spearman 高说明特征值和寿命位置有稳定排序关系。

\feature{mag__time__mean_abs} 与 \feature{mag__time__mean} 得分完全一致。

这说明它们在当前数据上高度冗余。

下游 baseline 可以二者都放入 full feature set，也可以在 compact subset 中二选一。

\feature{h__time__mean_abs}、\feature{h__time__rms}、\feature{h__time__std} 也在 RUL top 区域。

它们分别对应水平通道幅值、能量和离散度。

这些特征在 Step H 的三个工况内 top10 count 均较高。

不过 Step I 的跨工况检查显示，pure horizontal 特征比 magnitude 特征更容易受工况均值变化影响。

\begin{figure}[H]
\centering
\includegraphics[width=0.92\textwidth]{xjtu/main/rul_top_features.png}
\caption{XJTU-SY main split 下 RUL 任务的 top features。}
\label{fig:xjtu_rul_top_features}
\end{figure}

图 \ref{fig:xjtu_rul_top_features} 读法如下。

首先看 top 特征的 family。

可以看到前列集中在 magnitude/horizontal amplitude 和 dispersion。

这支持“退化趋势主要由振动幅值和能量上升表达”的结论。

其次要识别 Reference 特征。

\feature{mag__time__rms} 虽然排第一，但只能作为 label-source sanity reference。

第三要看非 label-source 的替代特征。

\feature{mag__time__mean}、\feature{mag__time__mean_abs} 和 \feature{mag__time__std} 提供了可以进入 baseline 的主候选。

因此 XJTU-SY RUL 最终推荐不是简单取 rank 1，而是取“高分且非 label-source 且稳定性检查支持”的 magnitude amplitude subset。

\subsection{Health State 特征分析}

Health State 的主划分结果同样由 amplitude family 驱动。

Step F 中 \feature{h__time__mean_abs} 的 health score 为 0.996，rank 1。

它的 mutual information 为 0.623，Fisher score 为 1.864，ANOVA F 为 4366.3。

这三个指标共同说明，不同伪健康阶段上的水平通道幅值分布差异很明显。

\feature{h__time__rms} 和 \feature{h__time__std} 分别 rank 2 和 rank 3。

它们几乎是同一物理现象的两个统计表达：水平振动能量和离散度随退化阶段增强。

如果只看 Step F，Health State 似乎应直接使用 horizontal amplitude family。

但是 Step I 提供了更严格的跨工况视角。

在 cross-condition ranking 中，vertical/magnitude features 排名更靠前，horizontal features 出现 downgrade。

原因不是 horizontal 特征无效，而是其均值和方差更容易随工况变化。

因此最终推荐把 XJTU-SY Health State 的 A 级特征放在 \feature{mag__time__mean}、\feature{mag__time__mean_abs}、\feature{mag__time__std}，把 horizontal family 放入 B 级。

\begin{figure}[H]
\centering
\includegraphics[width=0.92\textwidth]{xjtu/main/health_state_boxplots.png}
\caption{XJTU-SY main split 下 Health State 的箱线图。}
\label{fig:xjtu_health_boxplots}
\end{figure}

图 \ref{fig:xjtu_health_boxplots} 需要看三个层次。

第一，看 healthy 到 severe 的中位数是否有明显阶梯。

水平通道幅值类特征在主划分中通常呈现随阶段上升的趋势。

第二，看箱体是否分离。

若相邻阶段箱体高度重叠，说明该特征更适合趋势辅助，而不适合作为单独分类依据。

第三，看 Reference。

\feature{mag__time__rms} 分离强并不意外，因为它参与了 HI/FPT 伪标签构造。

所以它应作为 sanity reference，而不是 Health State 独立证据。

综合 Step F、Step H、Step I，XJTU-SY Health State 的结论是：

\begin{itemize}[leftmargin=2em]
  \item 主划分中 horizontal amplitude 很强；
  \item 工况内检查仍支持 amplitude family；
  \item 跨工况后 magnitude features 更稳；
  \item Reference 特征需要单独处理。
\end{itemize}

\subsection{Early Fault 特征分析}

Early Fault 是 XJTU-SY 中最敏感、最需要 caveat 的任务。

主划分 Step F 中 \feature{h__time__mean_abs}、\feature{h__time__std}、\feature{h__time__rms} 分别位于 Early Fault 前三。

\feature{h__time__mean_abs} 的 early fault score 为 0.996，AUC 为 0.949，Cohen's d 为 1.514。

这说明在主划分训练集上，FPT 前后水平通道幅值差异很明显。

但是 condition-wise 分析改变了最终解释。

Step H 显示：

\begin{itemize}[leftmargin=2em]
  \item Condition 1 中 spectral entropy 更突出；
  \item Condition 2 中 peak-to-peak / shock-like 特征更突出；
  \item Condition 3 中 horizontal amplitude 更突出。
\end{itemize}

Step I 的 cross-condition 分析进一步显示，主划分中的 horizontal Early Fault claim 需要降级。

这不是否定主划分，而是说明 Early Fault 特征更接近工况相关的早期变化模式。

\begin{figure}[H]
\centering
\includegraphics[width=0.92\textwidth]{xjtu/main/early_fault_effects.png}
\caption{XJTU-SY main split 下 Early Fault 的特征效应图。}
\label{fig:xjtu_early_fault_effects}
\end{figure}

图 \ref{fig:xjtu_early_fault_effects} 的读法是比较 FPT 前后分布。

如果 FPT 后均值明显上升，且分布重叠较小，该特征对 early fault detection 有用。

主划分图中 horizontal amplitude 的效应很强。

但该图只代表主划分训练证据。

结合 Step H/I 后，最终不能把 horizontal amplitude 写成 XJTU-SY 全局稳定 Early Fault 特征。

最终策略是：

\begin{itemize}[leftmargin=2em]
  \item \feature{mag__time__mean} 和 \feature{mag__time__mean_abs} 保留为 A 级主线；
  \item \feature{mag__time__std}、\feature{v__time__std}、\feature{v__time__mean_abs} 保留为 B 级；
  \item horizontal amplitude 与 spectral entropy 放入 C 级，作为条件性辅助；
  \item \feature{mag__time__rms} 只作 Reference。
\end{itemize}

\subsection{工况内稳定性分析}

工况内分析来自 Step H。

它把三种工况分开，各自做 leave-one-bearing-out。

这样可以避免一个工况的样本数量或幅值尺度支配总体 ranking。

\begin{table}[htbp]
\centering
\caption{XJTU-SY condition-wise 与 cross-condition 稳定性}
\label{tab:xjtu_condition_stability}
\scriptsize
\begin{tabularx}{\textwidth}{lYYYY}
\toprule
任务 & C1 观察 & C2 观察 & C3 观察 & Cross-condition 决策 \\
\midrule
RUL & amplitude family 稳定 & amplitude family 稳定 & amplitude family 稳定 & \feature{mag__time__mean}, \feature{mag__time__mean_abs}, \feature{mag__time__std} 保持 A 级。\\
Health State & 水平幅值强 & horizontal RMS/std 更强 & 水平/幅值仍强 & pure horizontal 有降级；magnitude features 更适合 robust baseline。\\
Early Fault & spectral entropy 明显 & peak-to-peak/shock-like 特征明显 & horizontal amplitude 明显 & 最工况敏感；magnitude amplitude 保留主线，horizontal/spectral 放入 C 级。\\
Distribution shift & \multicolumn{3}{Y}{跨工况时 train/val/test 均值和方差变化明显，尤其 horizontal amplitude 在 test 均值低于 train。} & train-only ranking 可用，但解释必须带工况 caveat。\\
\bottomrule
\end{tabularx}
\end{table}


RUL 的稳定性最好。

三个工况中，amplitude family 都反复进入 top 区域。

这说明 RUL 退化趋势主要由振动幅值和能量变化表达，而不是完全依赖某一个工况。

Health State 的稳定性次之。

主 split 中 horizontal amplitude 强，但 C2/C3 与跨工况结果提示通道偏好会改变。

所以最终推荐中把 magnitude features 放到 A 级，horizontal features 放到 B 级。

Early Fault 最不稳定。

C1 中 spectral entropy 更强，C2 中 shock-like peak-to-peak 更强，C3 中 horizontal amplitude 更强。

这说明 Early Fault 更接近“具体工况下的早期变化检测”，而不是一个可以无 caveat 迁移的统一信号。

\subsection{跨工况鲁棒性分析}

Step I 使用 train=35Hz12kN，val=37.5Hz11kN，test=40Hz10kN。

该设置比主划分更严格。

ranking 仍然只在 train 上 fit。

val/test 主要用于分布和图像检查。

这意味着 Step I 的 top features 首先反映 35Hz12kN 的训练分布，但它能暴露“如果换工况，某些主划分结论是否危险”。

跨工况结果显示，RUL 中 \feature{v__time__mean_abs}、\feature{v__time__rms}、\feature{v__time__std} 和 magnitude features 仍然靠前。

Health State 中 magnitude features 保持稳定，而部分 horizontal features 被 downgrade。

Early Fault 中 \feature{v__spectral__entropy} 在 train 工况很强，但被标为 condition-sensitive，而不是全局推荐。

因此 Step I 的作用不是替代 Step F，而是给最终推荐加上等级：

\begin{itemize}[leftmargin=2em]
  \item 跨工况仍稳的进入 A/B；
  \item 主划分强但跨工况不稳的进入 C 或 downgrade；
  \item label-source 进入 Reference。
\end{itemize}

\begin{figure}[H]
\centering
\includegraphics[width=0.92\textwidth]{xjtu/cross/feature_recommendation_matrix.png}
\caption{XJTU-SY cross-condition 下的特征推荐矩阵。}
\label{fig:xjtu_cross_recommendation_matrix}
\end{figure}

图 \ref{fig:xjtu_cross_recommendation_matrix} 用于看跨工况 ranking 后哪些特征仍服务多个任务。

它支持 magnitude/vertical amplitude 在跨工况下更稳的判断。

它也提醒我们：Early Fault 的 spectral entropy 虽然显眼，但属于条件性信号。

\subsection{XJTU-SY 最终决策}

\begin{table}[htbp]
\centering
\caption{XJTU-SY 三任务推荐特征摘要}
\label{tab:xjtu_top_features}
\small
\begin{tabularx}{\textwidth}{lYYY}
\toprule
任务 & A 级 & B 级 & C 级或 Reference \\
\midrule
RUL & \feature{mag__time__mean}, \feature{mag__time__mean_abs}, \feature{mag__time__std} & \feature{h__time__mean_abs}, \feature{h__time__rms}, \feature{h__time__std}, \feature{v__time__mean_abs}, \feature{v__time__std} & Reference: \feature{mag__time__rms} \\
Health State & \feature{mag__time__mean}, \feature{mag__time__mean_abs}, \feature{mag__time__std} & \feature{h__time__mean_abs}, \feature{h__time__std}, \feature{h__time__rms} & C: \feature{v__spectral__entropy}; Reference: \feature{mag__time__rms} \\
Early Fault & \feature{mag__time__mean}, \feature{mag__time__mean_abs} & \feature{mag__time__std}, \feature{v__time__std}, \feature{v__time__mean_abs} & C: \feature{h__time__mean_abs}, \feature{h__time__std}, \feature{h__time__rms}, \feature{v__spectral__entropy}; Reference: \feature{mag__time__rms} \\
\bottomrule
\end{tabularx}
\end{table}


XJTU-SY 的主线特征集为 \config{manual_basic}。

RUL 使用 magnitude amplitude/dispersion 作为主线。

Health State 使用 magnitude features 作为稳健主线，horizontal features 作为 secondary。

Early Fault 使用 magnitude amplitude 作为安全主线，同时保留 horizontal/spectral C 级特征用于 ablation。

\config{manual_tsfresh_basic} 因 full-size tsfresh 内存压力被 blocked and deferred。

后续若要重新评估 tsfresh，必须先做 memory-aware backend，再评价它是否提供超出 amplitude/statistics 的新增信息。
